\section{Specification and Design}
\label{sec:design}

\subsection{The board decides, the client draws}
Responsibility is split so that the board is the only component that can decide
anything. The STM32 holds the game state, validates every request and emits the
single authoritative reply. The PC client renders what it is told and forwards
clicks; it enforces no rule, which is why a modified or hostile client cannot
produce an illegal position.

Table~\ref{tab:split} sets out who owns what. The pattern to notice is that
every row the client owns is presentation, and every row that could change the
outcome of the game sits on the board.

\begin{table}[htbp]
\caption{Where each decision is taken. Board modules are under \texttt{server/}, client modules under \texttt{client/}.}
\label{tab:split}
\centering
\footnotesize
\begin{tabular}{p{3.7cm}p{1.2cm}p{1.9cm}}
\hline
\textbf{Decision} & \textbf{Owner} & \textbf{Module} \\
\hline
Board contents & board & \texttt{game.c} \\
Whose turn it is & board & \texttt{game.c} \\
Whether a move is legal & board & \texttt{game.c} \\
Win and draw detection & board & \texttt{game.c} \\
Framing and dispatch & board & \texttt{protocol.c} \\
Button and LED meaning & board & \texttt{app.c} \\
When to redraw & client & \texttt{gui.py} \\
What a reply means to the user & client & \texttt{session.py} \\
\hline
\end{tabular}
\end{table}

The firmware is layered so that only one file is hardware-dependent.
\texttt{server/core/} is plain C with no HAL. \texttt{server/stm32/app.c} owns
the button, the LED and the main-loop behaviour but still has no HAL in it, and
\texttt{server/stm32/port\_stm32.c} is the single file that touches registers.
That is what allows the host simulator in \texttt{sim/} to compile the identical
core and application sources.

\subsection{A rejected move returns before the board is written}
Fig.~\ref{fig:fsm} gives the server FSM together with the move-validation
branch. The three states are IDLE, IN\_PROGRESS and GAME\_OVER. IDLE accepts
\texttt{CONNECT} and \texttt{NEW}; IN\_PROGRESS accepts \texttt{MOVE},
\texttt{STATE} and \texttt{RESIGN}; GAME\_OVER accepts \texttt{NEW},
\texttt{STATE} and \texttt{DISCONNECT}. Any command a state does not accept is
answered \texttt{NACK:NOGAME} and changes nothing, which is why those edges are
not drawn individually.

The validation branch is the part worth reading closely. A \texttt{MOVE} passes
three checks in a fixed order, range then turn then occupancy, and each one
exits to its own \texttt{NACK} before the board is touched. Only when all three
pass is the cell written, and only then is the win and draw test run. The
ordering is deliberate: a client that names the wrong player is told so, rather
than being told about a cell it had no business addressing.

\begin{figure*}[t]
\centering
\begin{tikzpicture}[>=Stealth, font=\small,
    st/.style={draw, rounded corners=2pt, minimum width=2.4cm, minimum height=0.9cm},
    chk/.style={draw, diamond, aspect=2.2, inner sep=1pt, font=\footnotesize},
    act/.style={draw, minimum width=2.2cm, minimum height=0.8cm, font=\footnotesize},
    nack/.style={draw, dashed, minimum height=0.7cm, font=\footnotesize},
    note/.style={font=\footnotesize, align=left},
    lbl/.style={font=\scriptsize, align=center}]

  % ---------------- panel (a): the three states
  \node[note, anchor=west, font=\footnotesize\bfseries] at (-1.0,1.9) {(a) Server states};

  \node[st] (idle) at (0.6,0)   {IDLE};
  \node[st] (prog) at (6.2,0)   {IN\_PROGRESS};
  \node[st] (over) at (12.4,0)  {GAME\_OVER};

  \draw[->] (idle) to[loop above] node[lbl, above]{CONNECT / OK} (idle);
  \draw[->] (prog) to[loop above] node[lbl, above]{STATE / BOARD} (prog);
  \draw[->] (idle) -- node[lbl, above]{NEW /\\OK:BOARD} (prog);
  \draw[->] (prog) -- node[lbl, above]{WIN, DRAW, RESIGN /\\WIN:p, DRAW, OK:GAMEOVER} (over);
  \draw[->] (over) to[out=140, in=40, looseness=0.7]
        node[lbl, above]{NEW / OK:BOARD} (prog);
  \draw[->] (over.south) -- ++(0,-0.8) coordinate (d)
        -- node[lbl, below]{DISCONNECT / OK} (d -| idle.south) -- (idle.south);

  \node[note, anchor=north west, text width=15.5cm] at (-1.0,-2.0)
    {Any command a state does not accept is answered \texttt{NACK:NOGAME} and
     changes nothing, so those edges are not drawn.};

  % ---------------- panel (b): what MOVE does inside IN_PROGRESS
  \node[note, anchor=west, font=\footnotesize\bfseries] at (-1.0,-3.1)
    {(b) The \texttt{MOVE} validation branch, taken inside IN\_PROGRESS};

  \node[act] (mv) at (0.6,-4.4) {\texttt{MOVE:}$i$};
  \node[chk] (c1) at (4.3,-4.4) {$0 \le i \le 8$?};
  \node[chk] (c2) at (8.0,-4.4) {player = turn?};
  \node[chk] (c3) at (11.7,-4.4) {cell empty?};
  \node[act] (wr) at (15.0,-4.4) {write the cell};

  \draw[->] (mv) -- (c1);
  \draw[->] (c1) -- node[lbl, above]{yes} (c2);
  \draw[->] (c2) -- node[lbl, above]{yes} (c3);
  \draw[->] (c3) -- node[lbl, above]{yes} (wr);

  \node[nack] (n1) at (4.3,-6.0) {\texttt{NACK:RANGE}};
  \node[nack] (n2) at (8.0,-6.0) {\texttt{NACK:NOTURN}};
  \node[nack] (n3) at (11.7,-6.0) {\texttt{NACK:OCCUPIED}};
  \draw[->] (c1) -- node[lbl, right]{no} (n1);
  \draw[->] (c2) -- node[lbl, right]{no} (n2);
  \draw[->] (c3) -- node[lbl, right]{no} (n3);

  \node[note, anchor=north, text width=3.4cm] at (15.0,-5.3)
    {then test the eight lines:\\
     win or draw $\rightarrow$ GAME\_OVER,\\
     otherwise \texttt{OK:BOARD} and
     the turn swaps};

  \node[note, anchor=north west, text width=12.0cm] at (-1.0,-7.0)
    {The three checks run in this order and each returns before the cell is
     written, so every dashed exit leaves the board, the turn and the status
     exactly as they were.};
\end{tikzpicture}
\caption{Server FSM (a) and the \texttt{MOVE} validation branch (b). The
validation chain sits inside the IN\_PROGRESS state: all three rejections exit
above the single write to the board.}
\label{fig:fsm}
\end{figure*}

\subsection{The client never blocks on the board}
Fig.~\ref{fig:gui} is the client flow. On start it opens the serial port and
sends \texttt{CONNECT}; the New Game button sends \texttt{NEW}; a click on cell
$i$ sends \texttt{MOVE:}$i$ and then repaints from whatever board string comes
back. A \texttt{NACK} is shown verbatim in the status line and the board is left
untouched, because the board the client holds is only ever a copy of the
server's. After a win or a draw the cells are disabled until a new game is
started, and \texttt{STATE} is issued once more so the final position is drawn
from the server rather than from the client's own guess.

The window does not freeze while a reply is outstanding. A daemon reader thread
blocks on the transport and pushes each line it receives onto a
\texttt{queue.Queue}; the Tk main loop drains that queue every 50\,ms through
\texttt{after()} and does all drawing there. A board that never answers
therefore costs a background thread, not the user interface, and the reader
applies its own timeout so a dead link surfaces as a message rather than a hang.

\begin{figure}[htbp]
\centering
\begin{tikzpicture}[>=Stealth, node distance=0.72cm, font=\footnotesize,
    bx/.style={draw, rounded corners=2pt, text width=3.1cm, align=center,
               minimum height=0.6cm},
    df/.style={draw, diamond, aspect=2.4, inner sep=1pt, align=center},
    lbl/.style={font=\scriptsize, align=center}]

  \node[bx] (start) {open the port, send \texttt{CONNECT}};
  \node[bx, below=of start] (wait) {wait for input};
  \node[bx, below=of wait] (send) {send \texttt{NEW} or \texttt{MOVE:}$i$};
  \node[df, below=of send] (rep) {reply?};
  \node[bx, below=1.15cm of rep] (draw) {redraw the board from the reply};
  \node[bx, below=of draw] (over) {announce, disable the cells, send \texttt{STATE}};

  \draw[->] (start) -- (wait);
  \draw[->] (wait) -- (send);
  \draw[->] (send) -- (rep);
  \draw[->] (rep) -- node[lbl, right]{\texttt{OK:BOARD}} (draw);
  \draw[->] (draw) -- node[lbl, right]{\texttt{WIN} or \texttt{DRAW}} (over);

  % right, inner: a NACK goes straight back to waiting
  \draw[->] (rep.east) -- ++(0.62,0) |- node[lbl, right, pos=0.42]
    {\texttt{NACK}:\\no redraw} (wait.east);
  % right, outer: a finished game waits for a new one
  \draw[->] (over.east) -- ++(1.5,0) |- (wait.east);
  % left, inner: an ordinary redraw returns to waiting
  \draw[->] (draw.west) -- ++(-0.62,0) |- (wait.west);
  % left, outer: a timeout or a closed port
  \draw[->] (rep.west) -- ++(-1.5,0) |- node[lbl, left, pos=0.42]
    {timeout or\\closed port} (wait.west);
\end{tikzpicture}
\caption{Client flow. Every redraw is driven by a server reply; a
\texttt{NACK} and a timeout both return to waiting without touching the board.}
\label{fig:gui}
\end{figure}

\subsection{One USB cable carries both the flashing and the link}
Fig.~\ref{fig:circuit} shows the connection. USART1 is PA9 for transmit and
PA10 for receive, and on this board both are wired on-board to the ST-LINK,
which also presents a USB virtual COM port. One cable therefore does both jobs
and no external USB-to-UART converter is involved, which is what the course
confirmed on 5 October 2026. The link runs at 115\,200 baud, 8 data bits, no
parity, one stop bit, no flow control, and the PC side appears as
\texttt{/dev/ttyACM0} on Linux.

Button0 is on PA0 behind a pull-up, so a press reads low, and the application
debounces it over 20\,ms and then feeds a synthetic \texttt{NEW} through the
server rather than resetting the game behind the protocol's back. LED0 is on
PB0 and reports whose turn it is: solid while X is to move, off while O is, and
blinking at 400\,ms whenever no game is running, which covers both the idle
board before the first game and a finished one. A glance at the board therefore
says whose turn it is, or that there is no turn to take, without looking at the
GUI.

\begin{figure}[htbp]
\centering
\begin{circuitikz}[american, scale=0.74, transform shape,
                   every node/.append style={font=\footnotesize}]
  \node[draw, minimum width=1.2cm, minimum height=0.9cm] (pc)  at (0,0)    {PC};
  \node[draw, minimum width=2.1cm, minimum height=0.9cm] (stl) at (2.9,0)  {ST-LINK VCP};
  \node[draw, minimum width=2.0cm, minimum height=2.6cm] (mcu) at (7.2,-0.7) {STM32F051};

  \draw (pc.east) -- node[above, font=\scriptsize]{USB} (stl.west);

  % USART1 to the on-board ST-LINK
  \draw[->] ([yshift=0.2cm]stl.east) -- node[above, font=\scriptsize]{PA9 TX}
        ([yshift=0.75cm]mcu.west);
  \draw[<-] ([yshift=-0.2cm]stl.east) -- node[below, font=\scriptsize]{PA10 RX}
        ([yshift=0.25cm]mcu.west);

  % Button0 on PA0, active low
  \draw ([yshift=-0.7cm]mcu.west) -- ++(-1.5,0) coordinate (b)
        node[above, pos=0.55, font=\scriptsize]{PA0};
  \draw (b) to[push button, l_={\scriptsize Button0}] ++(0,-1.3) node[ground]{};

  % LED0 on PB0
  \draw ([yshift=0.5cm]mcu.east) -- ++(1.2,0) coordinate (d)
        node[above, pos=0.55, font=\scriptsize]{PB0};
  \draw (d) to[R, l^={\scriptsize $R_s$}] ++(0,-1.1)
        to[leD, l^={\scriptsize LED0}] ++(0,-1.1) node[ground]{};
\end{circuitikz}
\caption{UART path and board I/O. One USB cable carries both the programming
interface and the virtual COM port that USART1 is wired to. Button0 is active
low and starts a new game; LED0 is solid while X is to move, off while O is,
and blinks at 400\,ms whenever no game is running.}
\label{fig:circuit}
\end{figure}
